\section{从零搭一个 x86-64 教学子集：先能手算，再谈 OS}

前面讲了从电路到 OS 的大方向，但真正理解硬件，不能只读概念。我们在这一节搭一个极小但完整的 x86-64 教学子集。它不追求真实 x86-64 的全部复杂度，而是保留所有关键骨架：bit、字节、机器字、内存地址、RIP、RSP、通用寄存器、Intel 语法指令、取指译码执行、栈、中断、Ring 0/3 和地址保护。只要这台小机器能讲清，后面再看 Linux x86-64，就只是把同一组机制放大。

\begin{keyidea}
先不要问“Linux 怎么做”。先问：如果只有一块内存、几个寄存器和一组指令，程序如何运行？程序如何等待设备？程序出错如何停住？多个程序如何互不破坏？这些问题推出来的答案，就是 OS 的硬件根。
\end{keyidea}

\subsection{第一步：bit、byte、word 到底是什么}

硬件里最小的稳定信息单位是 bit，它只有两个状态：0 或 1。8 个 bit 组成 byte。CPU 通常一次更喜欢处理固定宽度的一组 bit，例如 16 位、32 位或 64 位，这组 bit 叫 word。word 不是数学概念，而是机器一次自然处理的数据宽度。

\begin{longtable}{p{0.18\linewidth}p{0.26\linewidth}p{0.44\linewidth}}
\toprule
单位 & 例子 & 含义 \\
\midrule
bit & \code{0} 或 \code{1} & 一个开关状态、一个标志位、权限位 \\
byte & \code{01000001} & 最小寻址单位常是 1 byte，可表示字符或二进制片段 \\
word & 16/32/64 bit & CPU 算术、地址和寄存器常用宽度 \\
cache line & 常见 64 byte 量级 & cache 一次搬运的数据块 \\
page & 常见 4KiB 起 & MMU 和 OS 管理内存的基本单位 \\
\bottomrule
\end{longtable}

同一串 bit 可以被解释成不同含义。比如 \code{01000001} 可以是无符号整数 65，也可以是 ASCII 字符 \code{A}，也可以是某条指令的一部分。机器只保存 bit；含义来自解释规则。OS 做 ELF 装载、页权限、系统调用参数检查，本质上都在维护这些 bit 的合法解释方式。

\subsection{第二步：数值表示为什么需要补码}

无符号整数最简单：n 个 bit 表示 \code{0} 到 \code{2^n-1}。8 位无符号数 \code{11111111} 是 255。问题是负数怎么表示。现代 CPU 通常用二进制补码：最高位有负权重。8 位补码中，\code{11111111} 表示 -1，\code{10000000} 表示 -128，\code{01111111} 表示 127。

\begin{lstlisting}
8-bit unsigned:
  00000000 = 0
  00000001 = 1
  11111111 = 255

8-bit two's complement:
  00000000 = 0
  00000001 = 1
  01111111 = 127
  10000000 = -128
  11111111 = -1
\end{lstlisting}

补码的好处是加法电路统一。硬件不需要为有符号加法重做一套加法器。比如 \code{1 + (-1)}：

\begin{lstlisting}
  00000001
+ 11111111
=1 00000000   // 舍弃第 9 位进位，结果是 0
\end{lstlisting}

OS 为什么关心整数表示？因为长度、地址、偏移、权限位和计数器都可能溢出。系统调用里 \code{offset + length} 若不检查溢出，就可能绕过边界检查；页号移位若用错有符号类型，就会产生错误地址。底层表示不是数学细节，而是安全边界。

\subsection{第三步：内存是一张按地址编号的字节表}

先把内存想成一个数组：

\begin{lstlisting}
memory[0] = byte
memory[1] = byte
memory[2] = byte
...
\end{lstlisting}

地址就是数组下标。若机器有 16 位地址，则最多能直接编号 \code{2^16 = 65536} 个 byte，也就是 64KiB。真实 64 位机器虚拟地址空间远大得多，但第一层理解就是“地址选中某个 byte”。

\begin{longtable}{p{0.2\linewidth}p{0.34\linewidth}p{0.34\linewidth}}
\toprule
问题 & 朴素答案 & 后续复杂化 \\
\midrule
地址是什么 & 内存数组下标 & 虚拟地址、物理地址、DMA 地址会分开 \\
读一个 word 怎么办 & 读连续多个 byte 组合 & 端序、对齐、cache line、缺页 \\
写内存是什么 & 把 byte 放到地址位置 & 权限、dirty bit、写回 cache、DMA 一致性 \\
非法地址怎么办 & 简单机器可能乱读乱写 & MMU 触发异常，OS 决定修复或杀进程 \\
\bottomrule
\end{longtable}

多字节数值还涉及端序。小端机器把低有效字节放在低地址，大端机器相反。比如 32 位值 \code{0x12345678} 在小端内存中是：

\begin{lstlisting}
address A+0: 0x78
address A+1: 0x56
address A+2: 0x34
address A+3: 0x12
\end{lstlisting}

端序影响二进制文件、网络协议、设备寄存器和调试内存转储。OS 读 ELF header、网络栈解析报文、驱动读设备描述符时，都必须知道字段的字节顺序。

\subsection{第三步补充：一次内存访问要回答哪些问题}

在最初模型里，\code{memory[address]} 像数组访问。真实机器上，即使是最普通的 load，也不是直接把数组下标交给内存芯片。CPU 必须先生成有效地址，再检查地址形式、翻译虚拟地址、检查权限、访问 cache 或内存，最后在没有异常时写回寄存器。

\begin{longtable}{p{0.18\linewidth}p{0.36\linewidth}p{0.34\linewidth}}
\toprule
步骤 & X64S 简化说法 & 真实 x86-64 会继续做什么 \\
\midrule
地址生成 & \code{regs[base] + disp} & base/index/scale/disp 组合，检查 canonical address \\
权限判断 & 是否在允许范围 & 页表 U/S、R/W、NX、PKU/SMEP/SMAP 等权限 \\
地址翻译 & 教学内存直接使用地址 & TLB 命中或页表遍历，失败则 page fault \\
数据取得 & 读连续 byte 组合成 word & cache line、对齐、跨页、内存类型、MMIO 副作用 \\
写回提交 & 更新目标寄存器 & 必须保证异常前后架构状态精确 \\
\bottomrule
\end{longtable}

对齐和跨页是理解内存访问的第一个坑。假设页大小是 \code{0x1000}，一条 8 字节 load 从 \code{0x0ffc} 开始，它会读 \code{0x0ffc..0x1003}，跨过页边界。前 4 字节所在页可能存在，后 4 字节所在页可能不存在。硬件不能把半个值写进 \code{rax} 再报错；架构上要么完整完成，要么在写回前产生异常。

\begin{lstlisting}
load64 from 0x0ffc:
  bytes 0x0ffc..0x0fff are in page A
  bytes 0x1000..0x1003 are in page B
  if page B is not present:
    raise page fault before architectural writeback
\end{lstlisting}

这解释了为什么内核复制用户 buffer 不能只检查首地址。用户传入 \code{buf,len} 后，范围可能跨很多页；中间某页可能没有映射，某页可能只有读权限，某页可能在复制过程中被另一个线程 \code{munmap}。所以 \code{copy\_from\_user} 一类路径要能在每次 fault 后判断：能否修复，能否睡眠，是否应返回 \code{-EFAULT}。

\subsection{第三步补充：地址空间不是一块连续平板}

初学时把内存画成连续数组是必要简化，但进程地址空间不是一块实心平板。它是很多虚拟区间的集合：有代码段、只读数据、堆、文件映射、匿名映射、栈、guard page、vdso，也有大量没有映射的空洞。

\begin{lstlisting}
process virtual address sketch:
  text:        mapped, read+execute
  rodata:      mapped, read only
  data/heap:   mapped, read+write
  gap:         unmapped
  mmap file:   mapped, permissions by VMA
  guard page:  unmapped
  stack:       mapped, grows down
  kernel half: supervisor only
\end{lstlisting}

VMA 说明“这个范围应该是什么”，PTE 说明“这个页现在实际映射到哪里”。VMA 可以存在而 PTE 还不存在；这表示合法但尚未分配或尚未读入。PTE 可以暂时只读而 VMA 允许写；这表示写时复制。这个区分后面会反复出现：OS 不是只维护一个地址到物理页的 map，而是在维护“意图”和“当前实现状态”两层。

\subsection{第四步：先定义一个 x86-64 教学 CPU}

我们定义一个 x86-64 教学子集，名字叫 X64S。为了方便手算，它只开放少量指令和 64KiB 教学内存，但寄存器命名、栈语义、系统调用入口和中断返回都沿用 x86-64 主线。它有通用寄存器 \code{rax/rbx/rcx/rdx/rsi/rdi/rbp/rsp}，一个指令指针 \code{rip}，一个状态寄存器 \code{rflags}，一个当前特权级 \code{cpl}。

\begin{longtable}{p{0.16\linewidth}p{0.28\linewidth}p{0.44\linewidth}}
\toprule
状态 & 初始解释 & OS 关联 \\
\midrule
\code{rip} & 下一条指令地址 & 中断、异常和上下文切换必须保存它 \\
\code{rsp} & 当前栈顶地址 & 函数调用、系统调用、内核栈切换依赖它 \\
\code{rax/rdi/rsi/...} & 通用计算寄存器 & 返回值、系统调用号、参数、临时状态 \\
\code{rflags} & zero、carry、interrupt enable & 分支、算术结果、中断开关 \\
\code{cpl} & Ring 3 或 Ring 0 & 特权级的最小模型 \\
\bottomrule
\end{longtable}

这台机器是教学简化版，但它足够解释 OS。真实 x86-64 有更多寄存器、更复杂的 RFLAGS、控制寄存器、MSR、段/TSS/IST 和虚拟化状态；核心问题不变：机器状态必须能被保存、恢复、检查和保护。

\subsection{第五步：定义几条指令}

X64S 使用 x86-64 Intel 语法展示指令，但只选少量常用语义。真实 x86-64 编码有前缀、opcode、ModRM/SIB、位移和立即数；教学时先看语义。

\begin{longtable}{p{0.18\linewidth}p{0.28\linewidth}p{0.42\linewidth}}
\toprule
指令 & 语义 & 用途 \\
\midrule
\asmcode{mov reg, imm} & \code{reg = imm} & 把常量放进寄存器 \\
\asmcode{mov reg, [base]} & \code{reg = memory[base]} & 从内存读数据 \\
\asmcode{mov [base], reg} & \code{memory[base] = reg} & 向内存写数据 \\
\asmcode{add dst, src} & \code{dst = dst + src} & 算术计算 \\
\asmcode{sub dst, src} & \code{dst = dst - src} & 比较和减法 \\
\asmcode{cmp a, b} & 设置 zero/carry 等标志 & 条件分支前比较 \\
\asmcode{jmp label} & \code{rip = label} & 无条件跳转 \\
\asmcode{jz label} & zero 为 1 时跳转 & 条件分支 \\
\code{syscall} & 进入 Ring 0 系统调用入口 & 用户请求内核服务 \\
\code{iretq/sysretq} & 返回较低特权级或快速返回 & 受控返回用户态 \\
\bottomrule
\end{longtable}

真实机器会把 I/O 设计成端口 I/O 或 MMIO。本书主线使用 MMIO。教学重点是：指令不是魔法，它只是固定格式的 bit，译码器按 opcode 和寻址模式选择数据通路动作。

\subsection{第六步：取指、译码、执行的完整循环}

现在 X64S 能跑程序了。CPU 每次从 \code{rip} 指向的内存取指令字节，解码出指令长度和操作数，默认让 \code{rip} 前进到下一条指令；分支、异常和返回指令可以覆盖这个默认下一地址。

\begin{lstlisting}
cpu_step():
  inst = decode(memory, rip)
  next_rip = rip + inst.length
  switch opcode:
    MOV_REG_IMM: regs[dst] = imm
    MOV_LOAD:    regs[dst] = read64(memory, regs[base] + disp)
    MOV_STORE:   write64(memory, regs[base] + disp, regs[src])
    ADD:         regs[dst] = regs[dst] + regs[src]
    CMP:         rflags.zero = (regs[a] == regs[b])
    JZ:          if rflags.zero: next_rip = target
  rip = next_rip
\end{lstlisting}

真实 CPU 可能因为流水线和预测更复杂，但架构语义仍要保证程序按正确顺序推进。中断和异常会保存 \code{rip}，所以 fault 保存当前指令还是 interrupt 保存下一条指令，是架构必须定义清楚的契约。

\subsection{第七步：手跑一个程序}

我们写一个程序，把内存 \code{0x2000} 和 \code{0x2008} 的两个 64 位数相加，结果写到 \code{0x2010}。

\begin{lstlisting}
mov rbx, 0x2000
mov rax, qword ptr [rbx]       ; rax = memory[0x2000]
mov rbx, 0x2008
mov rcx, qword ptr [rbx]       ; rcx = memory[0x2008]
add rax, rcx
mov rbx, 0x2010
mov qword ptr [rbx], rax
hlt
\end{lstlisting}

假设初始内存：

\begin{lstlisting}
memory[0x2000] = 7
memory[0x2008] = 5
\end{lstlisting}

执行结束后，\code{memory[0x2010] = 12}。这就是所有程序运行的根：指令改变寄存器，寄存器形成地址，内存保存数据，\code{rip} 决定下一步。OS 再复杂，也不能违反这个根模型。

为了避免“看懂代码但不会手算”，下面把每条指令执行后的关键状态列出来。假设每条教学指令长度记作 1，初始 \code{rip=0}，未列出的寄存器保持不变。

\begin{longtable}{p{0.10\linewidth}p{0.28\linewidth}p{0.22\linewidth}p{0.28\linewidth}}
\toprule
步 & 指令 & 执行后寄存器 & 执行后内存变化 \\
\midrule
0 & 初始 & \code{rip=0} & \code{[0x2000]=7,[0x2008]=5} \\
1 & \asmcode{mov rbx, 0x2000} & \code{rbx=0x2000,rip=1} & 无 \\
2 & \asmcode{mov rax, qword ptr [rbx]} & \code{rax=7,rip=2} & 无 \\
3 & \asmcode{mov rbx, 0x2008} & \code{rbx=0x2008,rip=3} & 无 \\
4 & \asmcode{mov rcx, qword ptr [rbx]} & \code{rcx=5,rip=4} & 无 \\
5 & \asmcode{add rax, rcx} & \code{rax=12,rip=5} & 无 \\
6 & \asmcode{mov rbx, 0x2010} & \code{rbx=0x2010,rip=6} & 无 \\
7 & \asmcode{mov qword ptr [rbx], rax} & \code{rip=7} & \code{[0x2010]=12} \\
\bottomrule
\end{longtable}

这张表训练的是 OS 读源码时的基本能力：每一步都问状态在哪里、谁改了它、失败时能否回滚。系统调用、page fault、调度切换只是把这张表中的“寄存器、内存、下一条指令”扩展成更多对象。

\subsection{第七步补充：把指令编码成字节}

只讲汇编文本仍然少了一层。CPU 不读取 \asmcode{mov rax, qword ptr [rbx]} 这样的字符串，它读取内存里的字节。汇编器把文本翻译成机器码，链接器和装载器把这些字节放到可执行映射中，CPU 从 RIP 指向的地址取字节并译码。

为了能手算，X64S 使用固定长度教学编码。真实 x86-64 是变长编码，本书后面会讲 REX、opcode、ModRM、SIB、disp 和 imm；这里先抓住共同结构：opcode 决定操作，寄存器字段决定读写谁，位移和立即数字段提供常量。

\begin{longtable}{p{0.18\linewidth}p{0.26\linewidth}p{0.44\linewidth}}
\toprule
字段 & X64S 教学含义 & 真实 x86-64 对应 \\
\midrule
opcode & \code{MOV\_LOAD}、\code{ADD}、\code{SYSCALL} & opcode byte 或 opcode map \\
dst/src & 目标和源寄存器编号 & ModRM.reg、ModRM.r/m、REX 扩展位 \\
mode & 寄存器、内存、立即数寻址 & ModRM.mod、SIB、rip-relative \\
disp & 地址偏移 & 8/32 位 displacement \\
imm & 指令内常量 & immediate bytes，小端编码 \\
\bottomrule
\end{longtable}

\begin{lstlisting}
X64S teaching bytes for mov rax, qword ptr [rbx + 16]:
  opcode = MOV_LOAD64
  dst    = RAX
  base   = RBX
  disp   = 16

decode result:
  operation = load64
  target register = rax
  effective address = rbx + 16
\end{lstlisting}

这层编码很重要，因为异常和调试都要回到字节地址。RIP 指向的不是“第几行汇编”，而是某个机器码字节。断点调试器把指令首字节临时改成断点指令；异常处理器用保存的 RIP 找到 faulting instruction；动态链接器修改 GOT/PLT 相关地址，让后续 call 走到正确函数。操作系统管理的是字节、页和权限，不是源代码行。

\subsection{第七步补充：取指和数据访问有两条权限路径}

同样是访问内存，取指和读写数据的权限并不完全一样。取指要求页面可执行；load 要求可读；store 要求可写；用户态还要求 U/S 位允许用户访问。于是代码页通常是 read+execute，不可写；栈和堆通常是 read+write，不可执行。

\begin{longtable}{p{0.18\linewidth}p{0.28\linewidth}p{0.42\linewidth}}
\toprule
访问 & 检查的核心权限 & OS 用途 \\
\midrule
fetch instruction & present、user、execute & NX、W\^X、代码段保护 \\
load data & present、user、readable & 普通读、copy\_from\_user、文件映射 \\
store data & present、user、writable & 写堆栈、COW fault、脏页跟踪 \\
atomic RMW & read + write + 原子性 & 锁、引用计数、futex 快路径 \\
\bottomrule
\end{longtable}

这说明“某地址映射了”不是完整答案。一个页面可以存在但不可写，存在但不可执行，存在但只给 Ring 0 使用。后面讲 COW 时，父子进程共享同一物理页，但 PTE 暂时清 writable；第一次写触发 page fault，内核复制页面后再恢复写权限。后面讲安全时，栈不可执行和内核页不可用户访问也是同一类硬件权限检查。

\begin{lstlisting}
if fetch from stack page and page.executable == false:
  raise page fault with instruction-fetch bit

if store to cow page and page.writable == false:
  raise page fault with write bit
  kernel may copy page and retry same store
\end{lstlisting}

\subsection{第八步：为什么需要栈}

没有栈，也能写直线程序；一旦需要函数调用、局部变量和返回地址，就需要一个按后进先出管理的区域。我们把 \code{rsp} 作为栈指针，约定栈向低地址增长。

\begin{lstlisting}
push rax:
  rsp = rsp - 8
  memory[rsp] = rax

pop rax:
  rax = memory[rsp]
  rsp = rsp + 8

call addr:
  push next_rip
  rip = addr

ret:
  pop rip
\end{lstlisting}

栈解释了很多 OS 现象。每个线程为什么要有自己的栈？因为函数调用链、返回地址和局部变量不能混在一起。系统调用为什么要从用户栈切到内核栈？因为用户栈不可信，内核不能把自己的返回地址和局部状态放在用户可改的内存里。

\subsection{第九步：内存布局从哪里来}

程序不能随便把代码、数据和栈混在同一块区域。我们给 X64S 一个简单内存布局：

\begin{lstlisting}
0x0000 - 0x0fff  kernel code
0x1000 - 0x1fff  user program code
0x2000 - 0x2fff  user data
0x3000 - 0x3fff  user stack
0xf000 - 0xffff  device registers
\end{lstlisting}

这张表一开始只是约定。若所有代码同权，用户程序仍然可以写 \code{0x0000} 覆盖 kernel code，也可以写 \code{0xf000} 操作设备。要把“约定”变成“保证”，需要硬件保护。OS 的复杂性就是从这里开始的。

\subsection{第十步：轮询设备为什么会浪费 CPU}

假设 X64S 有一个键盘设备，MMIO 状态寄存器在 \code{0xf000}，数据寄存器在 \code{0xf008}。状态 bit0 表示是否有字符。

\begin{lstlisting}
keyboard_poll:
  mov rbx, 0xf000
loop:
  mov rax, qword ptr [rbx]        ; read status
  and rax, 1
  cmp rax, 0
  jz loop
  mov rbx, 0xf008
  mov rcx, qword ptr [rbx]        ; read character
\end{lstlisting}

这个程序功能上正确，但 CPU 在没有按键时一直空转。若系统只有一个任务，可以接受；若还要同时控制电机、刷新显示、处理网络，就不能让 CPU 卡在一个设备等待上。于是需要中断：设备有事件时打断 CPU，而不是 CPU 一直问设备。

\subsection{第十一步：中断如何进入 CPU}

给 X64S 增加中断机制。设备可以通过 APIC/MSI-X 路径设置 pending vector。CPU 在指令边界，如果 \code{rflags.IF = 1} 且有可投递中断，就保存当前返回状态，关闭同级中断，跳到 IDT 中登记的中断入口。

\begin{lstlisting}
after each instruction:
  if rflags.IF and pending_interrupt:
    kernel_stack.push(rflags)
    kernel_stack.push(cs)
    kernel_stack.push(rip)
    rflags.IF = 0
    cpl = 0
    rip = idt[vector].handler
\end{lstlisting}

中断处理函数读取设备数据，清 pending，恢复 \code{rip/rflags}，返回被打断的程序。

\begin{lstlisting}
irq_return:
  iretq restores rip, cs, rflags, rsp, ss when returning to Ring 3
\end{lstlisting}

这就是 trap frame 的雏形。真实 CPU 保存更多寄存器，也有更复杂的栈切换和优先级，但本质是一样的：硬件必须留下足够状态，让内核处理事件后能回到被打断的位置。

\subsection{第十一步补充：IDT、TSS 和内核栈的关系}

在 x86-64 上，异常和中断不是直接跳到任意地址。CPU 用 vector 查 IDT，IDT entry 指出 handler 地址、门类型和目标特权级。若事件从 Ring 3 进入 Ring 0，CPU 还要找到可信内核栈。这个栈信息来自 TSS 里的 RSP0 或 IST 配置。

\begin{lstlisting}
ring 3 interrupt entry:
  vector = pending IRQ or exception number
  gate = IDT[vector]
  if gate target is ring 0 and current CPL is ring 3:
    new_rsp = TSS.rsp0 or selected IST stack
    switch to new_rsp
    push old SS, old RSP, RFLAGS, CS, RIP
  else:
    push frame on current privileged stack
  RIP = gate.handler
  CPL = 0
\end{lstlisting}

为什么不能继续用用户栈？因为用户栈由用户程序控制。它可能未映射，可能接近 guard page，可能故意伪造内容。若内核把 trap frame 和局部变量压到用户栈，用户态就能破坏内核返回路径。因此每个线程必须有自己的内核栈：线程 A 在系统调用里睡眠时，它的内核调用链还在；线程 B 进入内核不能覆盖 A 的栈。

\begin{longtable}{p{0.18\linewidth}p{0.34\linewidth}p{0.36\linewidth}}
\toprule
结构 & 保存什么 & 错误配置后果 \\
\midrule
IDT & vector 到 handler 的入口描述 & page fault、timer IRQ 无法分发 \\
TSS.RSP0 & Ring 3 进入 Ring 0 时的栈 & trap frame 写到错误位置或三重故障 \\
IST & 特殊异常专用栈 & double fault、NMI 等无法可靠处理 \\
per-task kernel stack & 线程在内核里的调用链 & 睡眠线程被其他线程覆盖现场 \\
\bottomrule
\end{longtable}

这也是为什么“创建线程”不是只分配用户栈。内核还要准备内核栈、初始切换上下文、返回用户态所需的寄存器帧、页表归属和调度状态。没有这些，线程无法被中断、无法系统调用、无法在内核里阻塞后再恢复。

\subsection{第十二步：定时器为什么能让 OS 抢回 CPU}

中断不只来自键盘，也可以来自定时器。给 X64S 加一个每 1000 个周期触发一次的 local APIC timer IRQ。即使用户程序写了死循环，定时器仍会周期性把 CPU 带回内核。

\begin{lstlisting}
user_bad_program:
  jmp user_bad_program

timer_irq:
  save current registers into current_task.context
  choose next runnable task
  restore next_task.context
  iretq
\end{lstlisting}

这就是抢占式调度的硬件根。没有定时器，OS 只能相信程序主动 \code{yield}；有了定时器，OS 可以在程序不合作时收回 CPU。调度器不是凭空切换线程，它依赖定时器中断、上下文保存和恢复。

\subsection{第十三步：为什么需要用户态和内核态}

现在 X64S 有中断了，但用户程序仍然可以写设备寄存器、改中断向量、覆盖内核代码。我们加入 \code{cpl}：Ring 3 表示用户态，Ring 0 表示内核态。再规定：

\begin{itemize}
\tightlist
\item 只有 Ring 0 能访问 \code{0x0000-0x0fff} 的内核代码数据。
\item 只有 Ring 0 能访问 \code{0xf000-0xffff} 的设备 MMIO。
\item 只有 Ring 0 能修改 IDT、CR3、MSR 和保护配置。
\item Ring 3 执行特权指令会触发异常。
\end{itemize}

\begin{lstlisting}
memory_access(addr, type):
  if cpl == 3 and addr in kernel_or_device_range:
    raise protection_fault
  else:
    perform access
\end{lstlisting}

这一步把“不要乱写”变成“写不了”。操作系统最基本的安全性，不是靠应用自觉，而是靠 CPU 在每次敏感访问时强制检查。

\subsection{第十四步：系统调用是受控入口}

有了用户态/内核态，用户程序不能直接写设备，那它如何输出字符？答案是系统调用。X64S 使用 x86-64 的 \code{syscall} 指令。用户程序把系统调用号放进 \code{rax}，把参数放进 \code{rdi/rsi/rdx/r10/r8/r9}，执行 \code{syscall}；CPU 切到 Ring 0 入口；内核检查系统调用号和参数，再代表用户操作设备。

\begin{lstlisting}
user:
  mov rax, 1        ; syscall number: write_char
  mov rdi, 'A'      ; argument
  syscall

hardware:
  save user rip/rflags into rcx/r11
  cpl = 0
  rip = syscall_entry_msr

kernel:
  if rax == write_char:
    check permission
    write device register
  return to user
\end{lstlisting}

现在可以看清：系统调用不是函数调用。函数调用不改变 \code{cpl}，不能绕过用户权限；系统调用由硬件提供受控入口，只有入口处的内核代码能获得特权。

\subsection{第十五步：base/limit 是最小内存保护}

先不引入复杂页表，给 X64S 教学模型每个用户程序两个界限寄存器：\code{BASE} 和 \code{LIMIT}。用户态访存时，硬件检查地址是否落在范围内。真实 x86-64 不用这种简单 base/limit 管普通用户地址空间，而是依赖页表和 U/S 权限位；这里先用它说明“硬件必须参与隔离”。

\begin{lstlisting}
user_translate(addr):
  if addr >= LIMIT:
    raise protection_fault
  physical = BASE + addr
  return physical
\end{lstlisting}

这让每个程序以为自己从地址 0 开始，硬件把它平移到不同物理区域。程序 A 的 \code{addr=0x0100} 和程序 B 的 \code{addr=0x0100} 可以落到不同物理位置。它已经能实现基本隔离，但有明显局限：每个程序需要连续物理内存，不适合共享库，也不适合按需分配。

\subsection{第十五步补充：base/limit 为什么会被页表取代}

base/limit 的优点是硬件简单：一次比较和一次加法就能完成保护和重定位。但它把一个进程看成单个连续区间，这和真实程序的形态冲突。真实程序需要代码只读可执行、数据可写、栈向下增长、堆向上增长、共享库映射到任意位置、文件按需映射、guard page 捕获栈溢出。

\begin{longtable}{p{0.22\linewidth}p{0.34\linewidth}p{0.32\linewidth}}
\toprule
需求 & base/limit 的问题 & 分页的做法 \\
\midrule
稀疏地址空间 & 空洞也要占连续范围 & 未使用区间不分配页表下级或 PTE \\
不同权限 & 整段权限太粗 & 每页独立 R/W/X 和 U/S 权限 \\
共享库 & 难把同一物理内容映射到多进程 & 多个 PTE 指向同一物理页 \\
按需装载 & 必须提前准备连续物理内存 & present=0，fault 时再分配或读文件 \\
COW & 整段复制成本高 & PTE 只读共享，写 fault 时复制页 \\
\bottomrule
\end{longtable}

这并不是说 base/limit 没有价值。它是理解“硬件参与边界检查”的第一块台阶。页表只是把同一个思想细化到页粒度，并增加了稀疏性、共享、延迟分配和权限组合。现代 OS 的内存管理之所以强大，是因为硬件每次访问都愿意替内核检查这一页当前是否允许这种访问。

\subsection{第十六步：分页如何解决连续内存问题}

分页把虚拟地址切成页号和页内偏移。假设页大小 256 byte，虚拟地址高 8 位是虚拟页号，低 8 位是偏移。页表记录每个虚拟页映射到哪个物理页，以及权限。

\begin{lstlisting}
translate(va, access):
  vpn = va >> 8
  off = va & 0xff
  pte = page_table[vpn]
  if !pte.present:
    raise page_fault
  if access not allowed:
    raise protection_fault
  return (pte.ppn << 8) | off
\end{lstlisting}

分页解决了三件事。第一，虚拟连续不要求物理连续。第二，每页可以有独立权限。第三，不存在的页可以等访问时再分配或从文件读取。现代虚拟内存、COW、\code{mmap}、按需装载都从这里长出来。

\subsection{第十六步补充：四级页表如何节省空间}

真实 x86-64 常见 48 位 canonical 虚拟地址和 4KiB 页。4KiB 页内偏移需要 12 位，剩下的高位用来索引多级页表。四级页表把虚拟地址拆成 PML4、PDPT、PD、PT、offset，每级通常取 9 位，因为一个 4KiB 页表页可以容纳 512 个 8 字节 PTE。

\begin{lstlisting}
48-bit virtual address:
  bits 47..39  PML4 index
  bits 38..30  PDPT index
  bits 29..21  PD index
  bits 20..12  PT index
  bits 11..0   page offset
\end{lstlisting}

如果用单级线性页表，为整个低 48 位地址空间准备所有 PTE，会消耗巨大内存；多级页表只为实际用到的区间分配下级表。比如一个小进程只映射代码、堆、栈和少量共享库，大量虚拟空洞没有对应页表页。这就是“稀疏地址空间”的硬件数据结构。

\begin{lstlisting}
page_walk(va):
  pml4e = read(CR3 + pml4_index(va) * 8)
  if !pml4e.present: fault
  pdpte = read(pml4e.page + pdpt_index(va) * 8)
  if !pdpte.present: fault
  pde = read(pdpte.page + pd_index(va) * 8)
  if !pde.present: fault
  pte = read(pde.page + pt_index(va) * 8)
  if !pte.present or permission denied: fault
  return pte.frame + offset(va)
\end{lstlisting}

TLB 的意义也在这里变得具体。一次完整 page walk 可能读四级页表，任何一级还可能 cache miss。TLB 把最终翻译和权限缓存起来，让常见访问不必重复走表。但这也带来维护义务：PTE 改了，TLB 里的旧答案不会自动在所有 CPU 上消失。

\subsection{第十六步补充：PTE 位背后的 OS 决策}

PTE 不是单纯的物理页号。它是一份硬件和内核共同维护的小合同。硬件读取 present、writable、user、NX 等位来决定访问是否成立；硬件可能设置 accessed/dirty 位；内核读取这些位来做回收、写回、COW、安全策略和调试诊断。

\begin{longtable}{p{0.18\linewidth}p{0.34\linewidth}p{0.36\linewidth}}
\toprule
PTE 信息 & 硬件视角 & OS 用法 \\
\midrule
present & 这一页能否翻译到物理页 & demand paging、swap、file-backed lazy load \\
writable & store 是否允许 & COW、只读代码、共享映射保护 \\
user & Ring 3 是否能访问 & 用户/内核地址空间隔离 \\
NX & 是否禁止取指 & W\^X、栈不可执行、JIT 权限切换 \\
accessed & 页面是否被访问过 & LRU 近似、老化、回收候选选择 \\
dirty & 页面是否被写过 & writeback、匿名页回收、文件页落盘 \\
global/PCID 相关 & TLB 是否跨切换保留 & 降低上下文切换成本 \\
\bottomrule
\end{longtable}

这张表也解释了为什么 page fault handler 必须看 error code。一次 fault 可能是不存在页，也可能是写只读页，也可能是用户态访问 supervisor 页，也可能是取指 NX 页。处理方式完全不同：合法缺页可以分配，COW 写可以复制，权限违规要发信号或杀进程，内核访问坏用户指针可能要返回 \code{-EFAULT}。

\subsection{第十七步：page fault 不是一定错误}

如果用户访问一个 \code{present=0} 的页，硬件触发 page fault。内核查看这个地址是否属于合法 VMA。如果合法，例如栈增长或文件映射尚未读入，内核可以分配页、更新 PTE、返回用户态重试原指令。如果不合法，内核发送错误信号或杀死进程。

\begin{lstlisting}
page_fault(va, access):
  vma = find_vma(current, va)
  if vma == none:
    kill_process(SIGSEGV)
  else if access not allowed by vma:
    kill_process(SIGSEGV)
  else:
    page = allocate_or_load_page(vma, va)
    install_pte(va, page, permissions)
    return_to_faulting_instruction()
\end{lstlisting}

这一步非常重要。缺页不是“内存坏了”，而是硬件把某次无法完成的访问交给 OS 决策。虚拟内存的弹性来自这个机制。

\subsection{第十七步补充：缺页处理的决策树}

下面把 page fault 拆成更接近真实内核的决策树。读者要特别注意：硬件只告诉内核“这次访问不能按当前页表完成”，真正的语义判断由内核结合 VMA、PTE、当前上下文和访问类型完成。

\begin{lstlisting}
handle_page_fault(va, access, error_code, from_user):
  if !is_canonical(va):
    return signal_or_kernel_fault()

  vma = find_vma(mm, va)
  if vma == none:
    return signal_or_efault()

  if access violates vma.permissions:
    return signal_or_efault()

  pte = lookup_pte(mm, va)
  if pte.present and access == write and pte.cow:
    new_page = copy_old_page()
    install_writable_pte(va, new_page)
    flush_tlb_page(va)
    return retry

  if !pte.present and vma.anonymous:
    page = allocate_zero_page()
    install_pte(va, page, vma.permissions)
    return retry

  if !pte.present and vma.file_backed:
    page = page_cache_get_or_read(vma.file, va)
    install_pte(va, page, vma.permissions)
    return retry

  return signal_or_kernel_fault()
\end{lstlisting}

\begin{longtable}{p{0.22\linewidth}p{0.34\linewidth}p{0.32\linewidth}}
\toprule
fault 场景 & 内核判断依据 & 结果 \\
\midrule
栈增长 & fault 地址贴近栈 VMA 边界 & 扩展 VMA 或分配新页 \\
匿名懒分配 & VMA 允许访问但 PTE 不存在 & 分配零页并重试 \\
文件映射缺页 & VMA 指向文件页 & 从 page cache 或磁盘读入 \\
COW 写 & PTE 只读且标记 COW & 复制物理页，更新 PTE，重试 \\
非法用户地址 & 无 VMA 或权限不符 & 发送 SIGSEGV 或返回 \code{-EFAULT} \\
内核 bug & 内核态访问非法地址且不可恢复 & oops 或 panic \\
\bottomrule
\end{longtable}

“可恢复”不是由硬件单独决定的。用户态普通指令 fault，内核可以通过信号报告；内核在 \code{copy\_from\_user} 中 fault，可以通过异常表修正返回 \code{-EFAULT}；内核在任意位置解引用坏指针，通常就是内核 bug。相同的硬件 fault，在不同上下文下有不同软件语义。

\subsection{第十七步补充：异常、陷阱和中断的返回地址不同}

很多教材把 interrupt、trap、fault 都翻成“中断/异常”，读者容易混在一起。对 OS 来说，最关键的区别是保存的返回 RIP 指向哪里，以及内核返回后是否重试同一条指令。

\begin{longtable}{p{0.16\linewidth}p{0.30\linewidth}p{0.42\linewidth}}
\toprule
事件 & 保存的 RIP & 典型处理 \\
\midrule
fault & 指向出错指令 & page fault 修复后重试；不能修复则发信号 \\
trap & 指向下一条指令或定义好的后继位置 & 断点、单步调试，处理后继续 \\
interrupt & 指向下一条待执行指令边界 & 设备 IRQ、定时器 IRQ，处理后恢复原程序 \\
abort & 可能无法可靠恢复 & double fault、机器检查等严重错误 \\
\bottomrule
\end{longtable}

以 page fault 为例，\asmcode{mov rax, qword ptr [rdi]} 若访问的页不存在，但 VMA 允许读，内核装好 PTE 后返回，CPU 重新执行同一条 \asmcode{mov}。以定时器中断为例，被打断指令已经在边界上结束，返回后执行下一条。这个差异决定了内核保存 trap frame、更新统计和投递信号的方式。

\begin{lstlisting}
page fault:
  saved_rip = address of faulting load/store
  kernel may install mapping
  iretq retries saved_rip

timer interrupt:
  saved_rip = next instruction boundary
  kernel may schedule another task
  later iretq continues from saved_rip
\end{lstlisting}

这也是为什么“精确异常”如此重要。如果 CPU 不能保证 faulting RIP 和架构状态一致，内核就无法判断重试是否安全。

\subsection{第十八步：把所有状态放进 task}

如果要运行多个程序，OS 必须能暂停一个程序，恢复另一个程序。一个 task 至少保存这些状态：

\begin{lstlisting}
struct Task {
  gpr[16]
  rip
  rsp
  rflags
  cpl
  page_table
  kernel_stack
  state: Runnable | Running | Blocked | Zombie
}
\end{lstlisting}

上下文切换就是保存当前 task 的 CPU 状态，选择下一个 runnable task，把它的状态装回 CPU。真实 OS 还要保存浮点/SIMD、调度统计、信号、fd 表、凭据、namespace 等，但最小根就是这些。

\subsection{第十九步：从 X64S 教学子集映射到真实 x86-64}

现在可以把 X64S 的概念映射到真实机器。

\begin{longtable}{p{0.22\linewidth}p{0.34\linewidth}p{0.32\linewidth}}
\toprule
X64S 概念 & 真实 x86-64 对应 & OS 意义 \\
\midrule
\code{rip} & RIP & 异常返回和调试定位 \\
\code{rsp} & RSP、TSS/IST 栈切换 & 用户栈、内核栈、调度切换 \\
\code{cpl} & Ring 3/Ring 0 & 特权级和系统调用边界 \\
syscall entry & IA32\_LSTAR MSR、IDT gate & 受控进入内核 \\
page table & x86-64 四/五级页表、CR3、PCID & 地址空间和权限 \\
timer IRQ & local APIC timer、TSC deadline & 抢占和超时 \\
device range & PCIe BAR、ACPI resource & 驱动访问设备寄存器 \\
task context & \code{task\_struct/thread\_info/pt\_regs} 等 & 调度器保存恢复执行 \\
\bottomrule
\end{longtable}

真实机器复杂，但不是另一套逻辑。x86-64 的页表层级更多，权限位更多，系统调用入口更讲究，设备通过 PCIe 和 MSI-X，DMA 经过 IOMMU；但每一项都能在 X64S 的最小模型里找到原型。

\subsection{第十九步补充：x86-64 long mode 启动后有哪些关键控制点}

真实 x86-64 进入 64 位 long mode 后，OS 要配置一批“硬件查表入口”。这些入口不是普通全局变量，而是 CPU 在异常、系统调用、地址翻译和栈切换时直接读取的控制状态。

\begin{longtable}{p{0.16\linewidth}p{0.34\linewidth}p{0.38\linewidth}}
\toprule
控制点 & 保存什么 & OS 为什么必须初始化 \\
\midrule
CR3 & 当前顶级页表物理地址和 PCID 信息 & 每次地址翻译都从这里找到地址空间 \\
IDT & vector 到 interrupt/trap gate 的表 & page fault、timer IRQ、device IRQ 要有入口 \\
GDT/TSS & 段描述和特权栈信息 & 从 Ring 3 进 Ring 0 时选择内核栈，IST 处理特殊异常 \\
MSR LSTAR & \code{syscall} 的 64 位内核入口 & 用户态系统调用只能跳到受控入口 \\
MSR STAR/FMASK & syscall 段选择和 RFLAGS mask & 进入内核时清除不安全标志位 \\
APIC registers & 本地中断、定时器、IPI 控制 & 抢占、跨 CPU 调度、TLB shootdown \\
\bottomrule
\end{longtable}

如果 IDT 没设好，page fault 会变成无法处理的更严重异常；如果 TSS 的内核栈错误，用户态陷入内核时会把 trap frame 写到错误位置；如果 CR3 指向坏页表，取指和访存都会立即崩溃；如果 LSTAR 指向错误地址，\code{syscall} 就不是受控入口而是灾难。启动代码的价值，就是在允许用户程序运行前，把这些硬件入口全部变成可解释、可恢复、可审计的状态。

\begin{lstlisting}
minimal x86-64 kernel entry preparation:
  build early page tables
  load CR3
  load GDT and TSS
  install IDT entries for faults and IRQs
  write syscall MSRs
  initialize local APIC timer
  create first task and kernel stack
  only then return or jump to user mode
\end{lstlisting}

这张清单也解释了为什么 OS 实践卷不能一上来写“进程”。没有页表、异常入口、内核栈和定时器，进程只是结构体；有了这些硬件控制点，进程才变成可以被隔离、打断、恢复和销毁的执行实体。

\subsection{本节验收：能不能从头讲一遍机器运行}

\begin{rubricbox}
\begin{itemize}
\tightlist
\item 能解释 bit、byte、word、地址、内存数组之间的关系。
\item 能手跑一个 \code{mov/add/mov} 小程序，说明每步 \code{rip} 和寄存器如何变化。
\item 能解释为什么需要栈，以及 \code{call/ret} 如何使用栈保存返回地址。
\item 能解释轮询设备为什么浪费 CPU，中断如何保存 \code{rip/rflags} 并跳到处理函数。
\item 能解释定时器为什么让 OS 可以抢占死循环程序。
\item 能解释用户态为什么不能直接写设备寄存器，系统调用如何成为受控入口。
\item 能解释 base/limit 和分页分别解决什么问题，page fault 什么时候是可修复事件。
\item 能把 X64S 的 \code{rip/rsp/cpl/page table/timer IRQ} 映射到真实 x86-64 概念。
\end{itemize}
\end{rubricbox}
